\chapter{Ergebnisse}
\label{ch:ergebnisse}

\section{Esterbildung}
\Cref{fig:ester} zeigt den deutlichen Anstieg von Isoamylacetat mit der Temperatur.

\begin{figure}[htbp]
  \centering
  \begin{tikzpicture}
    \begin{axis}[
      width=0.85\linewidth, height=6.2cm,
      xlabel={Gärtemperatur [\si{\celsius}]}, ylabel={Konzentration [mg/L]},
      grid=major, grid style={black!10},
      legend pos=north west, legend style={draw=none},
      every axis plot/.append style={thick}
    ]
      \addplot[color=akzent, mark=*] coordinates {(16,1.1) (18,1.5) (20,1.9) (22,2.6) (24,3.1)};
      \addplot[color=orange, mark=square*] coordinates {(16,0.18) (18,0.22) (20,0.27) (22,0.31) (24,0.36)};
      \legend{Isoamylacetat (Banane), Ethylhexanoat (Apfel)}
    \end{axis}
  \end{tikzpicture}
  \caption{Esterkonzentration in Abhängigkeit der Gärtemperatur (Mittelwerte, $n=3$).}
  \label{fig:ester}
\end{figure}

\section{Verkostung}
Ab \SI{20}{\celsius} wurde der Unterschied von \SI{88}{\percent} der Personen
erkannt. Das bei \SI{18}{\celsius} vergorene Bier erhielt den besten
Gesamteindruck -- fruchtig, aber noch \enquote{sauber}.
